\section*{Chapitre \ref{ch:b3:facies-models} --- Modèles de faciès}

\begin{solution}{exo:b3:facies-models:1}
Du large vers la terre~: offshore, transition offshore avec sables de tempête,
\omterm{def:b3:facies-models:shoreface}{avant-côte}, plage, une plaine côtière à marais. Chaque milieu est venu se
superposer à son voisin situé vers le large~: le rivage s'est déplacé vers le large, c'est une régression. Une
transgression les aurait empilés dans l'autre sens, la plaine côtière à la
base et la vase offshore au sommet, souvent avec une surface d'érosion là où
l'\omterm{def:b3:facies-models:shoreface}{avant-côte} passait.
\end{solution}

\begin{solution}{exo:b3:facies-models:2}
(a) Un remplissage de chenal ou une \omterm{def:b3:facies-models:rivers}{barre de méandre}~: un affouillement graveleux à la base, des stratifications obliques,
puis des rides et une vase racinaire. (b) Un \omterm{def:b3:facies-models:shoreface}{avant-côte} progradant ou un front de delta~:
de la vase avec des bancs de sable de tempête et de la \omterm{def:b3:facies-models:hcs}{stratification oblique en mamelons} (HCS) en dessous,
un sable d'\omterm{def:b3:facies-models:shoreface}{avant-côte} à stratification oblique et une plage au-dessus. (c) Des sables
de \omterm{def:b3:facies-models:rivers}{rivière en tresses} empilés, une dune éolienne ou un haut-fond carbonaté~: des ensembles de stratifications obliques à taille de grain constante
avec de grands foresets de dunes dans le cas éolien, ou des
grains squelettiques et oolithiques dans le cas du haut-fond.
\end{solution}

\begin{solution}{exo:b3:facies-models:3}
$\phi = (2.65 - 2.40)/(2.65 - 1.0) = 0.152$~; avec de la saumure,
$0.25/1.55 = 0.161$.
\end{solution}

\begin{solution}{exo:b3:facies-models:4}
$e^{-2\pi z/\lambda} = e^{-\pi} = 0.043$ à $z = \lambda/2$~: environ
\qty{38}{m} et \qty{55}{m}. À \qty{20}{m}~: $e^{-2\pi\times20/75} = 0.19$
et $e^{-2\pi\times20/110} = 0.32$ du mouvement de surface.
\end{solution}

\begin{solution}{exo:b3:facies-models:5}
$(60 - 47.6)/(189 - 47.6) = 0.088$~; $(60 - 52.6)/(189 - 52.6) = 0.054$. Une
valeur de matrice incertaine de \qty{5}{\micro\second/ft} modifie la porosité
d'un tiers~: l'équation de moyenne temporelle exige une valeur de matrice vérifiée sur carottes.
\end{solution}

\begin{solution}{exo:b3:facies-models:6}
$(6.0 + 3.5 + 8.0 + 2.5)/40 = 20/40 = 0.50$. Volume poreux~: $20\times0.22
\times10^6 = \qty{4.4e6}{m^3}$ par kilomètre carré.
\end{solution}

\begin{solution}{exo:b3:facies-models:7}
Près du canyon~: Ta seulement, les divisions supérieures ayant été enlevées par l'écoulement suivant~; en descendant
l'éventail~: Tc (rides) et Td (silt) sous Te (vase). Le courant perd de la vitesse
et sa charge grossière avec la distance~; seule la fraction fine parcourt
\qty{30}{km}, et elle est déposée par un écoulement plus lent qui construit
des rides, et non par retombée d'une suspension dense.
\end{solution}

\begin{solution}{exo:b3:facies-models:8}
Une \omterm{def:b3:facies-models:rivers}{barre de méandre} d'une \omterm{def:b3:facies-models:rivers}{rivière à méandres}~: un résidu grossier sur le fond érodé du chenal,
du sable à stratification oblique du chenal profond, du sable à rides du sommet de la barre,
une vase de plaine d'inondation avec racines. La succession enregistre le chenal du fond à
la berge, si bien que le chenal avait environ \qty{9}{m} de profondeur (un peu moins, après
compaction de la vase). Pas à travers la base~: c'est une surface d'érosion, et
ce qui se trouvait en dessous a pu être n'importe quel dépôt plus ancien.
\end{solution}

\begin{solution}{exo:b3:facies-models:9}
$60/40\,000 = 0.0015$, \ang{0.09}~: une rampe. $300/1500 = 0.20$, \ang{11.3},
après un palier sommital de \qty{5}{m} de profondeur~: une \omterm{def:b3:facies-models:ramp}{plate-forme à barrière}, avec des grains de récif ou de haut-fond
en bordure, et des blocs et des écoulements granulaires sur la pente en contrebas.
\end{solution}

\begin{solution}{exo:b3:facies-models:10}
La ceinture met un temps $w/v$ pour passer en un point donné, pendant lequel le substrat
s'élève par rapport au niveau marin à la vitesse $A$ sous elle~: le sable remplit $h = A\,w/v$.
$2000\times0.2/20 = \qty{20}{m}$. Avec $A = 0$, l'\omterm{def:b3:facies-models:shoreface}{avant-côte} ne laisse que
l'épaisseur de son propre profil, un mince placage, puisque les ceintures glissent
les unes sur les autres sans s'accumuler. Le rapport $A/v$ est la tangente de l'angle de
la trajectoire énoncée dans la proposition du chapitre~\ref{ch:b3:sequence-stratigraphy}~:
l'épaisseur de chaque faciès est sa largeur multipliée par cette tangente.
\end{solution}

\begin{solution}{exo:b3:facies-models:11}
Densité~: $(2.65 - 2.30)/1.65 = 0.212$~; sonique~: $(86 - 52.6)/136.4 = 0.245$.
Avec $\phi = 0.245$, $\rho_{\mathrm{ma}} = (2.30 - 0.245)/(1 - 0.245) =
\qty{2.72}{g/cm^3}$~: plus dense que le quartz, comme pour un grès cimenté par
de la calcite (2.71) ou contenant des grains de calcite.
\end{solution}

\begin{solution}{exo:b3:facies-models:12}
De la vase offshore reposant sur un charbon met en contact des milieux non adjacents~: une
\omterm{def:b3:sequence-stratigraphy:parasequence}{surface d'inondation}, ici une \omterm{def:b3:sequence-stratigraphy:surfaces}{surface de transgression} d'érosion par les vagues, les débris de coquilles
formant un résidu. Au fur et à mesure que le rivage se déplaçait vers la terre, les vagues ont entaillé le fond
jusqu'à environ la \omterm{def:b3:facies-models:shoreface}{limite d'action des vagues de beau temps}~: la plage, la partie supérieure de l'\omterm{def:b3:facies-models:shoreface}{avant-côte} et
le sommet de la plaine côtière, sur environ \qty{10}{m} de coupe, peuvent
manquer.
\end{solution}

\begin{solution}{pb:b3:facies-models:1}
\textbf{1.} Craie (froide, faible densité), argile à partir de \qty{766}{m} (chaude), calcaire à partir de
\qty{866}{m} (froid, dense, rapide), grès à partir de \qty{1046}{m} (formes en
cloche), schiste à partir de \qty{1166}{m}.

\textbf{2.} Craie \qty{766}{m}, argile 100, calcaire 180, grès
\qty{120}{m}~; base de la couverture à $441.0 - 1166 = \qty{-725}{m}$.

\textbf{3.} $(2.71 - 2.03)/1.71 = 0.40$ près de la surface, $(2.71 -
2.36)/1.71 = 0.20$ à \qty{750}{m}~: la craie perd la moitié de sa porosité
sous le poids des roches de recouvrement, par compaction.

\textbf{4.} Ses grains sont des argiles de densité incertaine, et son eau est
en partie liée à elles~: les hypothèses de l'équation ne sont pas vérifiées.

\textbf{5.} Socle cristallin~: la couverture est traversée, et aucun réservoir n'est
attendu en dessous.

\textbf{6.} $(2.71 - 2.573)/1.71 = 0.080$.

\textbf{7.} $(58.9 - 47.6)/141.4 = 0.080$~: même résultat.

\textbf{8.} $(2.65 - 2.32)/1.65 = 0.200$.

\textbf{9.} $(79.9 - 52.6)/136.4 = 0.200$~: même résultat.

\textbf{10.} $0.33/1.55 = 0.213$, soit un tiers de plus.

\textbf{11.} Le grès, avec une porosité de \qty{20}{\percent} dans ses sables~;
les argilites au-dessus du calcaire constituent la roche couverture, et les boues de plaine d'inondation divisent le
grès en compartiments.

\textbf{12.} Une forme en cloche~: faible radioactivité gamma dans le sable à stratification oblique, qui remonte à travers
l'intervalle granodécroissant jusqu'aux fortes valeurs de la boue enracinée, sur une base nette.

\textbf{13.} Une \omterm{def:b3:facies-models:association}{succession granodécroissante}~: dépôt grossier résiduel au fond du chenal sur une base
érosive, dunes de chenal (stratifications obliques), sommet de la \omterm{def:b3:facies-models:rivers}{barre de méandre} (rides), plaine
d'inondation (boue enracinée).

\textbf{14.} Le chenal, sa \omterm{def:b3:facies-models:rivers}{barre de méandre} sur la berge interne, des levées et une
plaine d'inondation pourvue de sols, côte à côte.

\textbf{15.} Environ l'épaisseur du sable et de l'intervalle granodécroissant, $10.4 + 3.2 =
\qty{13.6}{m}$, un peu plus avant la compaction.

\textbf{16.} Une \omterm{def:b3:facies-models:rivers}{rivière à méandres} dont la ceinture de chenaux s'est déplacée à plusieurs reprises
par \omterm{def:b3:facies-models:rivers}{avulsion} pendant que le bassin subsidait, chaque nouveau chenal entaillant la plaine
d'inondation laissée par le précédent.

\textbf{17.} Non~: chaque base est érosive~; la boue de plaine située sous elle et le sable
de chenal situé au-dessus n'ont pas nécessairement été voisins.

\textbf{18.} $(60 - 30)/(105 - 30) = 0.40$.

\textbf{19.} Cas minimal~: $14\times(0.60 + 0.40\times0.25) = \qty{9.80}{m}$~;
cas maximal~: $16\times(0.65 + 0.40\times0.20) = \qty{11.68}{m}$.

\textbf{20.} Sable propre \qty{58.8}{m}~; intervalles granodécroissants \qty{33.4}{m}~; boue
\qty{27.8}{m}.

\textbf{21.} Épaisseur nette $= 58.8 + 0.40\times33.4 = \qty{72.2}{m}$~; \omterm{def:b3:facies-models:ntg}{rapport net sur
brut} $72.2/120 = 0.60$.

\textbf{22.} Fractions 0.20 et 0.60~: \qty{65.5}{m} et \qty{78.8}{m}, un
écart de \qty{13}{m} dû au seul choix du seuil de coupure.

\textbf{23.} $72.2\times0.20\times10^6 = \qty{1.44e7}{m^3}$ par
kilomètre carré.

\textbf{24.} L'épaisseur de sable net du grès triasique au puits Varden-2 est
\textbf{\qty[text-series-to-math=true]{72.2}{m}}.
\end{solution}
